\section{What Sets the Behavioral Rate?}
\label{sec:behavior-prediction}

The rate-distortion view suggests measuring, before training, what the data ask of the receiver. Within one receiver this works: the log-determinant of the correction-gradient spectrum (Appendix~\ref{app:correction-statistics}) ranks $R^\star(0.90)$ across 24 arms and 7 task families with Spearman $\rho=0.92$ on Mistral-7B and $0.59$ on Qwen2.5-7B (Appendix Figure~\ref{fig:predictor}a; selection-adjusted $p=0.0003$). It does not transfer. Locked before each new receiver's rates were measured, it scores $\rho=-0.1$ on Llama-3.1-8B, $-0.4$ on Qwen3-8B and $+0.1$ on a fresh Qwen2.5-7B panel.

The reason is that the rate belongs to the dataset--receiver pair, not to the dataset. Training the same three corpora on seven receivers (Appendix Figure~\ref{fig:predictor}b), XBRL tagging is the cheapest corpus on six of them but not on Qwen2.5-7B, where math is cheaper; and Qwen2.5-Math-7B needs 1.03 bits per value to keep 90\% of its math gain, against 0.41 for Qwen2.5-7B. No measure of the corpus alone, or of the corpus scored by the frozen model, orders two receivers on the same corpus better than chance (best receiver-pair sign accuracy 0.56).

One measure does. The share of the trained update's energy in its leading singular direction, fixed before a confirmation panel of 64 adapters (4 receivers, 34 corpora) was scored, orders receiver pairs correctly 73\% of the time on held-out receivers, against 48\% for the best correction statistic, and lowers held-out error from 0.167 (corpus only) to 0.129 bits per value. Adding the base model and the data, by measuring how much training-data gain each direction carries, fails the same test (17\%). The best predictor we have thus reads the trained adapter, and the same leading direction that carries the damage in Section~\ref{sec:recovery}; predicting the rate from the data and the receiver alone remains open.
